\clearpage
\VerdierPhysicalPageAnchor{0077}%
\VerdierSemanticAnchor{ch1-2-4-3-simple-complex-morphisms}
したがって、単純複体へ移ることにより、二つの関手の射

\VerdierSemanticAnchor{ch1-2-4-3-eq2}
\begin{equation}
  \begin{aligned}
    \operatorname{comp}_{\Pi}h\colon
      \operatorname{comp}_{\Pi}F
      &\longrightarrow \operatorname{comp}_{\Pi}F',\\
    \operatorname{comp}_{\Sigma}h\colon
      \operatorname{comp}_{\Sigma}F
      &\longrightarrow \operatorname{comp}_{\Sigma}F'
  \end{aligned}
  \tag{2.4.3.2}
\end{equation}
を得る。

\VerdierSemanticAnchor{ch1-2-4-3-composition}
$h$ と $h'$ が多加法的関手の合成可能な二つの射ならば、次が成り立つ：
\begin{equation*}
  \begin{aligned}
    \operatorname{comp}_{\Pi}hh'
      &=\operatorname{comp}_{\Pi}h\,\operatorname{comp}_{\Pi}h',\\
    \operatorname{comp}_{\Sigma}hh'
      &=\operatorname{comp}_{\Sigma}h\,\operatorname{comp}_{\Sigma}h'.
  \end{aligned}
\end{equation*}

\VerdierSemanticAnchor{ch1-2-4-3-sign-rule-ii}
射 $\operatorname{comp}_{\Pi}h$（それぞれ、
$\operatorname{comp}_{\Sigma}h$）は、符号規則 II (1.6.6) に従う
$\mathbb Z[n]$-関手の射である。

\VerdierSemanticAnchor{ch1-2-4-4}
\noindent\textbf{2.4.4.}\quad
関手 $F$ がいくつかの変数について反変、すなわち
$i\in J\subset[n]$ に対して $C_i=(C_i')^\circ$ であるとき、
(2.4.1) で行った同一視のもとで、関手 $\operatorname{comp}_{\Pi}F$ と
$\operatorname{comp}_{\Sigma}F$ は対応する変数について反変である。

\VerdierSemanticAnchor{ch1-2-5}
\subsection{ホモトピー}

\VerdierSemanticAnchor{ch1-2-5-1}
\noindent\textbf{2.5.1.}\quad
$X$ と $Y$ を $C$ の二つの $n$-重複体、
$f,g\colon X\rightrightarrows Y$ を二つの複体の射とする。
$f$ と $g$ を結ぶ\emph{ホモトピー}とは、射の族

\VerdierSemanticAnchor{ch1-2-5-1-family}
\begin{equation*}
  a=\bigl(a^{\sigma,i}\colon X^{\sigma+e_i}\longrightarrow Y^\sigma,
    \ \sigma\in\mathbb Z[n],\ i\in[n]\bigr)
\end{equation*}
で、次を満たすものをいう：

\VerdierSemanticAnchor{ch1-2-5-1-eq1}
\begin{equation}
  \begin{aligned}
    a^{\sigma+e_i-e_j,j}d_X^{\sigma,i}
      &=d_Y^{\sigma-e_j,i}a^{\sigma-e_j,j},
      &&\forall\sigma\in\mathbb Z[n],\
        \ \forall i,j\in[n],\ i\ne j,\\
    g^\sigma-f^\sigma
      &=\sum_{i\in[n]}\Bigl(
          a^{\sigma,i}d_X^{\sigma,i}
          +d_Y^{\sigma-e_i,i}a^{\sigma-e_i,i}\Bigr),
      &&\forall\sigma\in\mathbb Z[n].
  \end{aligned}
  \tag{2.5.1.1}
\end{equation}

\VerdierSemanticAnchor{ch1-2-5-1-notation}
$f$ と $g$ を結ぶホモトピー $a$ を $a\colon f\rightsquigarrow g$ と記す。

\VerdierSemanticAnchor{ch1-2-5-2}
\noindent\textbf{2.5.2.}\quad
$f,g,h\colon X\longrightarrow Y$ を三つの射、
$a\colon f\rightsquigarrow g$ と $b\colon g\rightsquigarrow h$ を
二つのホモトピーとする。ホモトピー
$b+a\colon f\rightsquigarrow h$ を

\VerdierSemanticAnchor{ch1-2-5-2-sum}
\begin{equation*}
  (b+a)^{\sigma,i}=b^{\sigma,i}+a^{\sigma,i},
  \qquad \sigma\in\mathbb Z[n],\ i\in[n]
\end{equation*}
で定める。

\VerdierSemanticAnchor{ch1-2-5-3}
\noindent\textbf{2.5.3.}\quad
$f,g\colon X\rightrightarrows Y$ を二つの複体の射、
$a\colon f\rightsquigarrow g$ をホモトピーとし、
$u\colon Y\longrightarrow Y'$（それぞれ、$v\colon X'\longrightarrow X$）
を複体の射とする。ホモトピー
$u\star a\colon uf\rightsquigarrow ug$（それぞれ、
$a\star v\colon fv\rightsquigarrow gv$）を

\VerdierSemanticAnchor{ch1-2-5-3-whiskering}
\begin{equation*}
  (u\star a)^{\sigma,i}=u^\sigma a^{\sigma,i}
  \qquad
  \text{（それぞれ、}(a\star v)^{\sigma,i}
    =a^{\sigma,i}v^{\sigma+e_i}\text{）}
\end{equation*}
で定める。
